\documentclass[11pt]{article}

\usepackage[a4paper,top=24mm,bottom=24mm,left=23mm,right=23mm,headheight=18pt,headsep=8mm]{geometry}
\usepackage{fontspec}
\usepackage{amsmath}
\usepackage{amssymb}
\usepackage{bm}
\usepackage{xcolor}
\usepackage{graphicx}
\usepackage{booktabs}
\usepackage{array}
\usepackage{tabularx}
\usepackage{longtable}
\usepackage{framed}
\usepackage{multicol}
\usepackage{hyperref}

\defaultfontfeatures{Ligatures=TeX}
\setmainfont{Noto Serif}
\setsansfont{Noto Sans}

\definecolor{WechatGreen}{HTML}{07C160}
\definecolor{WechatDark}{HTML}{057A3E}
\definecolor{Ink}{HTML}{17202A}
\definecolor{Muted}{HTML}{5B6673}
\definecolor{RuleGray}{HTML}{D9DEE5}
\definecolor{PaperGray}{HTML}{F5F7F9}
\definecolor{GreenTint}{HTML}{EAF8F0}
\definecolor{Blue}{HTML}{2F6FEB}
\definecolor{BlueTint}{HTML}{EEF5FF}
\definecolor{Orange}{HTML}{C76B00}
\definecolor{OrangeTint}{HTML}{FFF4E5}
\definecolor{Red}{HTML}{B42318}
\definecolor{RedTint}{HTML}{FFF0EE}
\definecolor{Purple}{HTML}{7357B7}
\definecolor{PurpleTint}{HTML}{F4F0FF}

\hypersetup{
  colorlinks=true,
  hypertexnames=false,
  linkcolor=WechatDark,
  citecolor=WechatDark,
  urlcolor=Blue,
  pdftitle={From Softmax Attention to Gated DeltaNet},
  pdfsubject={A first-principles reconstruction of Scientific Spaces archive 11823},
  pdfkeywords={softmax attention, Gated DeltaNet, delta rule, linear attention}
}

\setlength{\parindent}{0pt}
\setlength{\parskip}{0.58em}
\setlength{\fboxsep}{7pt}
\renewcommand{\arraystretch}{1.25}
\linespread{1.06}
\setcounter{tocdepth}{2}

\newcolumntype{Y}{>{\raggedright\arraybackslash}X}
\newcolumntype{C}{>{\centering\arraybackslash}X}
\newcolumntype{P}[1]{>{\raggedright\arraybackslash}p{#1}}

\newcommand{\vect}[1]{\bm{#1}}
\newcommand{\trans}{^{\mathsf T}}
\newcommand{\R}{\mathbb{R}}
\newcommand{\softmax}{\operatorname{softmax}}
\newcommand{\bigO}{\mathcal{O}}
\newcommand{\sourceeq}[1]{\hfill{\sffamily\scriptsize\color{Muted}SOURCE EQ. #1}}

\newcommand{\badge}[2]{%
  \begingroup\setlength{\fboxsep}{3.5pt}%
  \colorbox{#1}{\textcolor{white}{\sffamily\bfseries\scriptsize #2}}%
  \endgroup
}

\newenvironment{greenbox}[1]{%
  \def\FrameCommand{\fcolorbox{WechatGreen}{GreenTint}}%
  \MakeFramed{\advance\hsize-\width\FrameRestore}%
  {\sffamily\bfseries\color{WechatDark}#1}\par\smallskip
}{\endMakeFramed}

\newenvironment{bluebox}[1]{%
  \def\FrameCommand{\fcolorbox{Blue}{BlueTint}}%
  \MakeFramed{\advance\hsize-\width\FrameRestore}%
  {\sffamily\bfseries\color{Blue}#1}\par\smallskip
}{\endMakeFramed}

\newenvironment{orangebox}[1]{%
  \def\FrameCommand{\fcolorbox{Orange}{OrangeTint}}%
  \MakeFramed{\advance\hsize-\width\FrameRestore}%
  {\sffamily\bfseries\color{Orange}#1}\par\smallskip
}{\endMakeFramed}

\newenvironment{redbox}[1]{%
  \def\FrameCommand{\fcolorbox{Red}{RedTint}}%
  \MakeFramed{\advance\hsize-\width\FrameRestore}%
  {\sffamily\bfseries\color{Red}#1}\par\smallskip
}{\endMakeFramed}

\newenvironment{graybox}[1]{%
  \def\FrameCommand{\fcolorbox{RuleGray}{PaperGray}}%
  \MakeFramed{\advance\hsize-\width\FrameRestore}%
  {\sffamily\bfseries\color{Ink}#1}\par\smallskip
}{\endMakeFramed}

\newcommand{\flowbox}[3]{%
  \fcolorbox{#1}{#2}{%
    \parbox{0.75\linewidth}{\centering\sffamily\bfseries\color{Ink}#3}%
  }%
}

\makeatletter
\renewcommand\section{\@startsection{section}{1}{0pt}%
  {2.7ex plus 0.7ex minus 0.2ex}{1.1ex}%
  {\normalfont\sffamily\Large\bfseries\color{Ink}}}
\renewcommand\subsection{\@startsection{subsection}{2}{0pt}%
  {2.1ex plus 0.5ex minus 0.2ex}{0.8ex}%
  {\normalfont\sffamily\large\bfseries\color{Ink}}}
\renewcommand\subsubsection{\@startsection{subsubsection}{3}{0pt}%
  {1.7ex plus 0.4ex minus 0.2ex}{0.6ex}%
  {\normalfont\sffamily\normalsize\bfseries\color{WechatDark}}}

\def\ps@wechat{%
  \def\@oddhead{\vbox{%
    \hbox to\textwidth{{\sffamily\footnotesize\color{Muted}SOFTMAX ATTENTION $\boldsymbol{\rightarrow}$ GATED DELTANET}\hfil{\sffamily\footnotesize\color{Muted}\thepage}}%
    \vskip3pt{\color{WechatGreen}\hrule height 0.65pt}}}%
  \def\@evenhead{\@oddhead}%
  \def\@oddfoot{\hfil{\sffamily\scriptsize\color{Muted}Scientific Spaces 11823 --- English technical reconstruction}\hfil}%
  \def\@evenfoot{\@oddfoot}%
}
\makeatother
\pagestyle{wechat}

\begin{document}

\begin{titlepage}
  \thispagestyle{empty}
  \vspace*{13mm}
  {\color{WechatGreen}\rule{34mm}{2.2mm}}\par
  \vspace{13mm}
  {\sffamily\bfseries\fontsize{27}{32}\selectfont\color{Ink}
    From Softmax Attention\\[2mm]
    to Gated DeltaNet\par}
  \vspace{5mm}
  {\sffamily\fontsize{15}{20}\selectfont\color{Muted}
    A derivation from first principles: motivation, mechanism, mathematics,
    approximations, and limits\par}

  \vspace{15mm}
  \begin{greenbox}{THE IDEA IN ONE SENTENCE}
    Softmax attention already performs an exact prediction-error correction
    for one fixed query. A local approximation of the newest softmax weight,
    followed by a heuristic query-independent anchor, turns that correction
    into a shared matrix update with Gated DeltaNet's retain--erase--write
    structure.
  \end{greenbox}

  \vspace{8mm}
  \begin{tabularx}{\textwidth}{@{}>{\sffamily\bfseries\color{Ink}}p{34mm}Y@{}}
    Primary source & Jianlin Su, Scientific Spaces archive 11823 \cite{su2026linearizing} \\
    Reference model & Gated Delta Networks, arXiv:2412.06464 \cite{yang2024gated} \\
    Format & English, equation-first, editorial explainer \\
    Claim boundary & Analytic reconstruction; no experiment or benchmark claim \\
  \end{tabularx}

  \vfill
  {\sffamily\small\color{Muted}Prepared for the Autoresearch Ideas repository}\par
  {\sffamily\small\color{Muted}23 July 2026}\par
\end{titlepage}

\tableofcontents
\clearpage

\section{The whole argument on one page}

\begin{center}
  \flowbox{Blue}{BlueTint}{Causal softmax attention for one fixed query}\par
  {\color{Muted}$\Downarrow$ \quad \badge{Blue}{EXACT REARRANGEMENT}}\par
  \flowbox{Blue}{BlueTint}{Old output $+$ newest weight $\times$ prediction error}\par
  {\color{Muted}$\Downarrow$ \quad \badge{Orange}{DIAGONAL TAYLOR APPROXIMATION}}\par
  \flowbox{Orange}{OrangeTint}{A query-dependent recurrent correction}\par
  {\color{Muted}$\Downarrow$ \quad \badge{Purple}{AFFINE STATE ANSATZ}}\par
  \flowbox{Purple}{PurpleTint}{A quadratic-in-query term prevents exact closure}\par
  {\color{Muted}$\Downarrow$ \quad \badge{Orange}{HEURISTIC ANCHOR}}\par
  \flowbox{Orange}{OrangeTint}{Replace the residual query by $k_t-\bar{k}_t$}\par
  {\color{Muted}$\Downarrow$ \quad \badge{WechatGreen}{STRUCTURAL MATCH}}\par
  \flowbox{WechatGreen}{GreenTint}{One shared state with retain $+$ erase $+$ write}\par
\end{center}

\begin{redbox}{READ THIS BEFORE THE EQUATIONS}
  The conclusion is not ``softmax attention equals Gated DeltaNet.'' The path
  contains a local Taylor expansion, an affine modeling assumption, and a
  norm-dependent heuristic choice. What emerges is a GDN-shaped approximation
  whose algebra reveals a meaningful relationship between attention and online
  error correction.
\end{redbox}

The derivation answers a practical question. Full causal attention can retrieve
from every prior token, but it forms a triangular score grid with
$1+2+\cdots+T$ query--key comparisons. A recurrent model updates one bounded
state per token. Can the behavior of softmax attention be compressed into such
a state without discarding the mechanism that corrects old associations? The
source article's answer is: approximately, and the correction naturally looks
like Gated DeltaNet.

\section{Foundations: the objects and operations}

\subsection{Dimensions and roles}

We use one attention head and a causal prefix ending at time $t$:

\[
  \vect q,\vect k_i\in\R^{d_k},\qquad
  \vect v_i,\vect o_t,\vect b_t\in\R^{d_v},\qquad
  \vect A_t\in\R^{d_v\times d_k}.
\]

\begin{tabularx}{\textwidth}{@{}p{27mm}p{31mm}Y@{}}
\toprule
\textbf{Object} & \textbf{Shape} & \textbf{Meaning} \\
\midrule
$\vect q$ & $d_k$-vector & A request: which stored pattern is relevant now? \\
$\vect k_i$ & $d_k$-vector & An address or feature direction for token $i$. \\
$\vect v_i$ & $d_v$-vector & The content returned if key $i$ is selected. \\
$\vect A_t$ & $d_v\times d_k$ matrix & A compressed map from query directions to value directions. \\
$\bar{\vect k}_t,\bar{\vect v}_t$ & vectors & Arithmetic prefix means through time $t$. \\
\bottomrule
\end{tabularx}

\subsection{Dot product: alignment becomes a scalar}

The score $\vect q\trans\vect k$ is a dot product:

\[
  \vect q\trans\vect k
  =\sum_{r=1}^{d_k}q_r k_r
  =\lVert\vect q\rVert\,\lVert\vect k\rVert\cos\theta.
\]

It is large and positive when the vectors point in similar directions, zero
when they are orthogonal, and negative when they oppose each other. Attention
uses this scalar as a compatibility score. The usual $1/\sqrt{d_k}$ scaling can
be absorbed into $\vect q$ or $\vect k$ for this derivation.

\subsection{Softmax: scores become normalized mixture weights}

For logits $s_i=\vect q\trans\vect k_i$,

\[
  \alpha_i=\frac{e^{s_i}}{\sum_j e^{s_j}},
  \qquad \alpha_i>0,
  \qquad \sum_i\alpha_i=1.
\]

Exponentiation rewards larger scores; normalization makes the result a convex
mixture. If the values are vectors, $\sum_i\alpha_i\vect v_i$ lies in their
weighted span.

\begin{graybox}{A SMALL NUMERICAL EXAMPLE}
  With logits $(-0.2,0,0.2)$, exact softmax gives approximately
  $(0.2693,0.3289,0.4018)$. The first-order approximation developed later gives
  $(0.2667,0.3333,0.4000)$. They are close because the logits are tightly
  clustered; the approximation need not remain close when they spread out.
\end{graybox}

\subsection{Outer product: a key--value association becomes a matrix}

The outer product $\vect v\vect k\trans$ has shape $d_v\times d_k$. Reading it
with a query gives

\[
  (\vect v\vect k\trans)\vect q
  =\vect v(\vect k\trans\vect q).
\]

The query first measures its alignment with the stored key, then scales the
stored value by that amount. This identity is the algebraic foundation of
matrix-state linear attention: a sum of outer products stores many approximate
associations in one matrix.

\section{Starting point: causal softmax attention}

For one fixed query $\vect q$, attention over the first $t$ key--value pairs is

\begin{equation}
  \vect o_t(\vect q)
  =\sum_{i=1}^{t}\alpha_{t,i}(\vect q)\vect v_i,
  \qquad
  \alpha_{t,i}(\vect q)
  =\frac{e^{\vect q\trans\vect k_i}}{Z_t(\vect q)},
  \qquad
  Z_t(\vect q)=\sum_{i=1}^{t}e^{\vect q\trans\vect k_i}.
  \tag{1}\label{eq:attention}
\end{equation}
\sourceeq{1}

\badge{Blue}{DEFINITION} Equation~\eqref{eq:attention} says that the query
scores every visible key, softmax normalizes those scores, and the resulting
weights mix the corresponding values. ``Causal'' means only positions
$1,\ldots,t$ are visible at position $t$.

\subsection{The direct first-order route gives vanilla linear attention}

Let $\bar{\vect k}_t=t^{-1}\sum_i\vect k_i$ and
$\bar{\vect v}_t=t^{-1}\sum_i\vect v_i$. A first-order expansion of every
softmax coefficient gives

\begin{align}
  \vect o_t(\vect q)
  &\approx \frac1t\sum_{i=1}^{t}
    \left[1+\vect q\trans(\vect k_i-\bar{\vect k}_t)\right]\vect v_i
    \notag\\
  &=\bar{\vect v}_t+
    \left(\frac1t\sum_{i=1}^{t}\vect v_i\vect k_i\trans
    -\bar{\vect v}_t\bar{\vect k}_t\trans\right)\vect q.
  \tag{2}\label{eq:vala}
\end{align}
\sourceeq{2}

Once the softmax approximation is made, the second line is exact algebra. The
matrix is a cross-covariance-like summary of keys and values. This is already a
linear-attention state, but it is essentially a Vanilla Linear Attention
(VaLA) variant. The source article asks whether a more targeted derivation can
retain the Delta Rule's ability to correct an existing association.

\section{The exact Delta-shaped recurrence}

Define the unnormalized numerator
$\vect N_t=\sum_{i=1}^{t}e^{\vect q\trans\vect k_i}\vect v_i$. Splitting off
the newest term gives

\begin{align}
  \vect o_t
  &=\frac{\vect N_t}{Z_t}
   =\frac{\vect N_{t-1}+e^{\vect q\trans\vect k_t}\vect v_t}{Z_t}
   =\frac{Z_{t-1}\vect o_{t-1}+e^{\vect q\trans\vect k_t}\vect v_t}{Z_t}.
  \tag{3}\label{eq:split}
\end{align}
\sourceeq{3}

Because $Z_t=Z_{t-1}+e^{\vect q\trans\vect k_t}$ and
$\alpha_{t,t}=e^{\vect q\trans\vect k_t}/Z_t$,

\begin{equation}
  \boxed{\vect o_t
  =\vect o_{t-1}
  +\alpha_{t,t}(\vect v_t-\vect o_{t-1})}.
  \tag{4}\label{eq:exactdelta}
\end{equation}
\sourceeq{4}

\begin{bluebox}{EXACT FOR ONE FIXED QUERY}
  Equation~\eqref{eq:exactdelta} has the canonical online-learning form
  \[
    \text{new prediction}
    =\text{old prediction}
    +\text{step size}\times\text{prediction error}.
  \]
  The error is $\vect v_t-\vect o_{t-1}$ and the adaptive step size is the
  newest attention weight $\alpha_{t,t}$. No approximation has appeared yet.
\end{bluebox}

\subsection{Why this exact RNN is still quadratic}

The suppressed dependence matters: the state is really
$\vect o_{t-1}(\vect q)$. In a Transformer, each position has a different
query $\vect q_1,\ldots,\vect q_T$. Query $\vect q_t$ needs a recurrence over
its entire prefix, so the amount of work is

\[
  1+2+\cdots+T=\frac{T(T+1)}2=\bigO(T^2).
\]

An actually linear recurrent model needs one state that is updated only from
keys and values, independent of whichever query will read it later. The rest of
the derivation constructs an approximation to that shared state.

\section{The local softmax approximation}

\subsection{First-order expansion from the ground up}

Let $s_i=\vect q\trans\vect k_i$, let
$\bar s=t^{-1}\sum_i s_i$, and define centered logits
$\delta_i=s_i-\bar s$. Then $\sum_i\delta_i=0$. For small $\delta_i$,

\[
  e^{s_i}=e^{\bar s}e^{\delta_i}
  =e^{\bar s}\left(1+\delta_i+\frac{\delta_i^2}{2}+\cdots\right).
\]

At first order the denominator is

\[
  \sum_j e^{s_j}
  \approx e^{\bar s}\sum_j(1+\delta_j)
  =t e^{\bar s},
\]

because the centered deviations cancel. Therefore

\begin{equation}
  \alpha_{t,i}(\vect q)
  \approx \frac1t(1+\delta_i)
  =\frac1t\left[1+\vect q\trans(\vect k_i-\bar{\vect k}_t)\right].
  \tag{T}\label{eq:taylor}
\end{equation}

This approximation still sums to one, but it no longer guarantees positivity.
Its omitted error begins at second order in the spread of the centered logits.
It is a local expansion around uniform attention, not a global replacement for
softmax \cite{su2026softmax}.

\subsection{Approximate only the newest coefficient}

Apply Equation~\eqref{eq:taylor} only to $\alpha_{t,t}$ in the exact recurrence.
Writing $\delta\vect k_t=\vect k_t-\bar{\vect k}_t$ gives

\begin{align}
  \vect o_t
  &\approx \vect o_{t-1}
    +\frac1t\left(1+\vect q\trans\delta\vect k_t\right)
      (\vect v_t-\vect o_{t-1}) \notag\\
  &=\left(1-\frac1t\right)\vect o_{t-1}+\frac1t\vect v_t
    +\frac1t(\vect v_t-\vect o_{t-1})
      \delta\vect k_t\trans\vect q.
  \tag{5}\label{eq:diagonal}
\end{align}
\sourceeq{5}

\badge{Orange}{LOCAL APPROXIMATION} This is a lighter approximation burden than
expanding all $\alpha_{t,i}$ at once, although it is not automatically accurate
for every sequence.

\subsection{Why the diagonal route can retain more structure}

For a fixed query, every old weight has the exact factorization

\begin{equation}
  \alpha_{t,i}
  =\alpha_{i,i}\prod_{j=i+1}^{t}(1-\alpha_{j,j}).
  \tag{6}\label{eq:factor}
\end{equation}
\sourceeq{6}

To see this, insert the telescoping normalizers:

\[
  \frac{e^{s_i}}{Z_t}
  =\frac{e^{s_i}}{Z_i}
   \frac{Z_i}{Z_{i+1}}\cdots\frac{Z_{t-1}}{Z_t},
  \qquad
  \frac{Z_{j-1}}{Z_j}=1-\frac{e^{s_j}}{Z_j}=1-\alpha_{j,j}.
\]

Approximating the local factors and multiplying them keeps cross-products such
as $(1+a)(1+b)=1+a+b+ab$. A single global first-order expansion would keep only
$1+a+b$. This motivates the diagonal strategy, but it is not a theorem that
the accumulated approximation error is uniformly smaller.

\section{A shared affine memory, and why it fails to close}

The smallest useful query-independent model is an affine readout:

\begin{equation}
  \widehat{\vect o}_t(\vect q)=\vect A_t\vect q+\vect b_t,
  \qquad \vect A_0=0,\quad\vect b_0=0.
  \tag{7}\label{eq:ansatz}
\end{equation}
\sourceeq{7}

\badge{Purple}{MODELING ANSATZ} The matrix and bias may depend on all observed
keys and values, but they must not depend on the query used at readout time.
This is deliberately an approximation: generic softmax attention is not affine
in $\vect q$.

Substitute Equation~\eqref{eq:ansatz} into Equation~\eqref{eq:diagonal}:

\begin{align}
  \vect A_t\vect q+\vect b_t
  \approx{}&\left(1-\frac1t\right)
    (\vect A_{t-1}\vect q+\vect b_{t-1})+\frac1t\vect v_t \notag\\
  &+\frac1t
    (\vect v_t-\vect b_{t-1}-\vect A_{t-1}\vect q)
    \delta\vect k_t\trans\vect q.
  \tag{8}\label{eq:substitution}
\end{align}
\sourceeq{8}

The query-free part closes cleanly:

\begin{equation}
  \vect b_t=\left(1-\frac1t\right)\vect b_{t-1}+\frac1t\vect v_t
  =\bar{\vect v}_t.
  \tag{9}\label{eq:bmean}
\end{equation}
\sourceeq{9}

But the apparent matrix update is

\begin{equation}
  \vect A_t=\left(1-\frac1t\right)\vect A_{t-1}
  +\frac1t(\vect v_t-\bar{\vect v}_{t-1}
  -\vect A_{t-1}\vect q)\delta\vect k_t\trans.
  \tag{10}\label{eq:notclosed}
\end{equation}
\sourceeq{10}

The visible $\vect q$ violates query independence. Equivalently,
Equation~\eqref{eq:substitution} contains

\[
  -(\vect A_{t-1}\vect q)(\delta\vect k_t\trans\vect q),
\]

which is quadratic in $\vect q$, while the left side is only affine. This is
the exact obstruction: first-order softmax does not imply first-order closure
of the recurrent state.

\section{The shortcut that only returns to VaLA}

One could simply drop $\vect A_{t-1}\vect q$ inside the residual, obtaining

\begin{equation}
  \vect A_t=\left(1-\frac1t\right)\vect A_{t-1}
  +\frac1t(\vect v_t-\bar{\vect v}_{t-1})
  (\vect k_t-\bar{\vect k}_t)\trans.
  \tag{11}\label{eq:drop}
\end{equation}
\sourceeq{11}

Because this is an incremental average, it unrolls exactly to

\begin{equation}
  \vect A_t=\frac1t\sum_{i=1}^{t}
  (\vect v_i-\bar{\vect v}_{i-1})
  (\vect k_i-\bar{\vect k}_i)\trans.
  \tag{12}\label{eq:unroll}
\end{equation}
\sourceeq{12}

This looks new, but two exact mean identities reveal that it is merely the
VaLA matrix in Equation~\eqref{eq:vala}. First,

\begin{align}
  t\bar{\vect v}_t\bar{\vect k}_t\trans
  -(t-1)\bar{\vect v}_{t-1}\bar{\vect k}_{t-1}\trans
  ={}&-\bar{\vect v}_{t-1}\bar{\vect k}_t\trans
  +\vect v_t\bar{\vect k}_t\trans
  +\bar{\vect v}_{t-1}\vect k_t\trans.
  \tag{13}\label{eq:meanidentity}
\end{align}
\sourceeq{13}

Summing this telescoping difference gives

\begin{equation}
  \sum_{i=1}^{t}\vect v_i\vect k_i\trans
  -t\bar{\vect v}_t\bar{\vect k}_t\trans
  =\sum_{i=1}^{t}
  (\vect v_i-\bar{\vect v}_{i-1})
  (\vect k_i-\bar{\vect k}_i)\trans.
  \tag{14}\label{eq:telescoping}
\end{equation}
\sourceeq{14}

\begin{graybox}{WHY THIS NEGATIVE RESULT MATTERS}
  Throwing away the troublesome prediction term also throws away the mechanism
  that erases an obsolete key--value association. The derivation has gone in a
  circle and recovered vanilla linear attention. To obtain a Delta-style state,
  the old prediction must be approximated more carefully rather than deleted.
\end{graybox}

\section{The error principle and the representative query}

Restore the query dependence explicitly in the exact fixed-query recurrence:

\begin{equation}
  \vect o_t(\vect q)=\vect o_{t-1}(\vect q)
  +\alpha_{t,t}(\vect q)
   [\vect v_t-\vect o_{t-1}(\vect q)].
  \tag{15}\label{eq:qexplicit}
\end{equation}
\sourceeq{15}

Suppose only the occurrence inside the residual is replaced by a representative
query $\vect c$. The introduced error is exactly

\begin{equation}
  \vect E_t(\vect q,\vect c)
  =\alpha_{t,t}(\vect q)
   [\vect o_{t-1}(\vect q)-\vect o_{t-1}(\vect c)].
  \tag{16}\label{eq:error}
\end{equation}
\sourceeq{16}

Choosing $\vect c=\vect q_t$ eliminates today's error at one query, but the
shared state must also serve unknown future queries. The article proposes a
minimax-inspired anchor: make the replacement exact at the query direction
where the newest token would receive its largest attention weight.

Under the linearized weight,

\[
  \alpha_{t,t}(\vect q)
  \approx \frac1t[1+\vect q\trans\delta\vect k_t].
\]

Without a norm constraint this has no finite maximum. If
$\lVert\vect q\rVert\le r$, Cauchy--Schwarz gives

\[
  \vect q\trans\delta\vect k_t
  \le \lVert\vect q\rVert\,\lVert\delta\vect k_t\rVert
  \le r\lVert\delta\vect k_t\rVert,
\]

with equality in the direction
$r\,\delta\vect k_t/\lVert\delta\vect k_t\rVert$. Assuming comparable query
and key norms, the source makes the simple magnitude choice

\begin{equation}
  \vect q^*=\delta\vect k_t=\vect k_t-\bar{\vect k}_t.
  \tag{17}\label{eq:qstar}
\end{equation}
\sourceeq{17}

\begin{orangebox}{WHAT IS RIGOROUS, AND WHAT IS HEURISTIC?}
  The maximizing \emph{direction} follows from Cauchy--Schwarz once a query-norm
  bound is imposed. The magnitude $\vect q^*=\delta\vect k_t$ additionally
  assumes comparable query/key scales. This does not solve a global minimax
  problem for Equation~\eqref{eq:error}; it anchors the approximation at one
  strategically selected direction.
\end{orangebox}

\section{The Gated DeltaNet-shaped update}

Define the centered key and value innovation

\[
  \vect x_t=\vect k_t-\bar{\vect k}_t,
  \qquad
  \vect y_t=\vect v_t-\bar{\vect v}_{t-1},
  \qquad
  g_t=\frac1t,
  \qquad
  \lambda_t=1-\frac1t.
\]

Substituting $\vect q^*=\vect x_t$ into Equation~\eqref{eq:notclosed} yields

\begin{align}
  \vect A_t
  &=\lambda_t\vect A_{t-1}
    +g_t(\vect y_t-\vect A_{t-1}\vect x_t)\vect x_t\trans
    \tag{18a}\label{eq:deriveddelta}\\
  &=\vect A_{t-1}
    (\lambda_t\vect I-g_t\vect x_t\vect x_t\trans)
    +g_t\vect y_t\vect x_t\trans.
    \tag{18b}\label{eq:derivedexpanded}
\end{align}
\sourceeq{18}

\begin{center}
\begin{tabularx}{0.96\textwidth}{@{}C@{\hspace{3mm}}C@{\hspace{3mm}}C@{}}
\fcolorbox{WechatGreen}{GreenTint}{\parbox{0.25\textwidth}{\centering\sffamily\bfseries
  RETAIN\par\smallskip$\lambda_t\vect A_{t-1}$\par\smallskip
  \normalfont\small Keep a decayed copy of global memory.}}
&
\fcolorbox{Orange}{OrangeTint}{\parbox{0.25\textwidth}{\centering\sffamily\bfseries
  ERASE\par\smallskip$-g_t(\vect A_{t-1}\vect x_t)\vect x_t\trans$\par\smallskip
  \normalfont\small Remove the old prediction along $\vect x_t$.}}
&
\fcolorbox{Blue}{BlueTint}{\parbox{0.25\textwidth}{\centering\sffamily\bfseries
  WRITE\par\smallskip$g_t\vect y_t\vect x_t\trans$\par\smallskip
  \normalfont\small Store the new centered target at $\vect x_t$.}}
\end{tabularx}
\end{center}

\subsection{The fundamental Delta Rule from gradient descent}

Suppose a matrix $\vect S$ should map key $\vect k$ to target value $\vect v$.
Use squared prediction loss

\[
  L(\vect S)=\frac12\lVert\vect v-\vect S\vect k\rVert^2.
\]

The matrix gradient is

\[
  \nabla_{\vect S}L=(\vect S\vect k-\vect v)\vect k\trans.
\]

One gradient-descent step with learning rate $\beta$ is

\begin{align}
  \vect S^+
  &=\vect S-\beta(\vect S\vect k-\vect v)\vect k\trans \notag\\
  &=\vect S+\beta(\vect v-\vect S\vect k)\vect k\trans \notag\\
  &=\vect S(\vect I-\beta\vect k\vect k\trans)
    +\beta\vect v\vect k\trans.
  \tag{D}\label{eq:deltarule}
\end{align}

The term $\vect S\vect k$ is the old prediction; the residual
$\vect v-\vect S\vect k$ is what remains to be learned. Multiplication by
$\vect k\trans$ localizes the matrix correction to that key direction.
Equation~\eqref{eq:deriveddelta} is this targeted correction plus a global
retention factor $\lambda_t$.

\subsection{A concrete two-dimensional update}

Let

\[
  \vect A=
  \begin{bmatrix}2&0.5\\0&1\end{bmatrix},\qquad
  \vect x=\begin{bmatrix}1\\0\end{bmatrix},\qquad
  \vect y=\begin{bmatrix}3\\-1\end{bmatrix},\qquad
  \beta=0.25.
\]

The old prediction is $\vect A\vect x=(2,0)\trans$ and the residual is
$(1,-1)\trans$. A plain Delta step gives

\[
  \vect A^+
  =\vect A+0.25
    \begin{bmatrix}1\\-1\end{bmatrix}
    \begin{bmatrix}1&0\end{bmatrix}
  =\begin{bmatrix}2.25&0.5\\-0.25&1\end{bmatrix}.
\]

Only the first column changes because $\vect x$ points along the first input
axis. The updated prediction $(2.25,-0.25)\trans$ has moved one quarter of the
way from the old prediction toward the target. This is the geometric meaning of
the rank-one erase/write operation.

\section{Comparison with canonical Gated DeltaNet}

A common Gated DeltaNet convention is \cite{yang2024gated}

\begin{equation}
  \vect S_t
  =\vect S_{t-1}
    [\alpha_t(\vect I-\beta_t\vect k_t\vect k_t\trans)]
   +\beta_t\vect v_t\vect k_t\trans.
  \tag{19}\label{eq:gdn}
\end{equation}
\sourceeq{19}

Both updates have the same three mechanisms: global retention, targeted erasure
of the current key's old prediction, and a rank-one write. For $t>1$,
Equation~\eqref{eq:derivedexpanded} can even be placed into the convention of
Equation~\eqref{eq:gdn} by the algebraic reparameterization

\[
  \alpha_t=\frac{t-1}{t},\qquad
  \beta_t=\frac1{t-1},\qquad
  \vect k'_t=\vect x_t,\qquad
  \vect v'_t=\frac{t-1}{t}\vect y_t.
\]

Indeed, $\alpha_t\beta_t=1/t$ and
$\beta_t\vect v'_t=\vect y_t/t$. At $t=1$, the centered key is zero and the
state can be initialized directly rather than using the singular
reparameterization.

\begin{tabularx}{\textwidth}{@{}p{32mm}YY@{}}
\toprule
 & \textbf{Taylor-derived state} & \textbf{Trained Gated DeltaNet} \\
\midrule
Retention & Deterministic $1-1/t$ & Learned, data-dependent gate \\
Update rate & Deterministic $1/t$ & Learned, data-dependent gate \\
Key & Centered $\vect k_t-\bar{\vect k}_t$ & Learned key projection \\
Value target & $\vect v_t-\bar{\vect v}_{t-1}$ & Learned value projection \\
Status & Approximation to softmax & Independently trained architecture \\
Evidence here & Analytic structure & No benchmark transferred by this derivation \\
\bottomrule
\end{tabularx}

\begin{redbox}{THE PRECISE CLAIM}
  There is an exact algebraic mapping between the displayed update forms after
  centering and rescaling their variables. There is not an equality between
  ordinary softmax attention and a trained Gated DeltaNet model. The latter
  learns its representations and gates from data, whereas the derivation fixes
  them through Taylor coefficients and prefix means.
\end{redbox}

\section{Readout, algorithm, and complexity}

The approximate output for the actual query at time $t$ is

\begin{equation}
  \widehat{\vect o}_t=\vect A_t\vect q_t+\bar{\vect v}_t.
  \tag{20}\label{eq:readout}
\end{equation}
\sourceeq{20}

The hat is intentional: it keeps visible that this is the output of the
approximating recurrent model, not an exact identity with the original softmax
output.

\subsection{One causal scan}

\begin{longtable}{@{}P{10mm}P{50mm}P{78mm}@{}}
\toprule
\textbf{Step} & \textbf{Operation} & \textbf{Purpose} \\
\midrule
1 & $g_t=1/t$, $\lambda_t=1-g_t$ & Set deterministic write and retention rates. \\
2 & $\bar{\vect k}_t=\lambda_t\bar{\vect k}_{t-1}+g_t\vect k_t$ & Update the key mean. \\
3 & $\vect x_t=\vect k_t-\bar{\vect k}_t$ & Isolate what is new about the key. \\
4 & $\vect y_t=\vect v_t-\bar{\vect v}_{t-1}$ & Isolate the value innovation against the old baseline. \\
5 & $\vect A_t=\lambda_t\vect A_{t-1}+g_t(\vect y_t-\vect A_{t-1}\vect x_t)\vect x_t\trans$ & Retain globally, erase the old prediction, and write the residual. \\
6 & $\bar{\vect v}_t=\lambda_t\bar{\vect v}_{t-1}+g_t\vect v_t$ & Update the value baseline. \\
7 & $\widehat{\vect o}_t=\vect A_t\vect q_t+\bar{\vect v}_t$ & Read the shared state with today's query. \\
\bottomrule
\end{longtable}

At $t=1$, $\bar{\vect k}_1=\vect k_1$, so $\vect x_1=0$,
$\vect A_1=0$, and the output is $\bar{\vect v}_1=\vect v_1$, exactly as it
should be for a one-token prefix.

\subsection{Sequence scaling}

\begin{tabularx}{\textwidth}{@{}p{31mm}YY@{}}
\toprule
 & \textbf{Causal softmax attention} & \textbf{Derived recurrent approximation} \\
\midrule
Sequence state & All prior keys and values: $\bigO(T(d_k+d_v))$ & Matrix and two means: $\bigO(d_kd_v+d_k+d_v)$ \\
Per-token work & Grows with prefix length & $\bigO(d_kd_v)$ rank-one update/read \\
Full sequence & $\bigO(T^2d_k)$ score interactions, plus value mixing & $\bigO(Td_kd_v)$ \\
\shortstack[l]{Scaling\\statement} & Quadratic in $T$ & Linear in $T$ for fixed head dimensions \\
\bottomrule
\end{tabularx}

Linear asymptotic sequence cost does not guarantee faster wall-clock execution.
Kernel fusion, parallel training scans, memory traffic, head size, and hardware
utilization all matter. The source article provides no timing measurements.

\section{Assumption and claim ledger}

\begin{longtable}{@{}P{29mm}P{38mm}P{77mm}@{}}
\toprule
\textbf{Class} & \textbf{Where} & \textbf{Boundary} \\
\midrule
\badge{Blue}{EXACT} & Equations 3--4 & Exact recurrence for one fixed query and its prefix. \\
\badge{Blue}{EXACT} & Equation 6 & Exact factorization of old attention weights into diagonal survival factors. \\
\badge{Orange}{TAYLOR} & Equations 2 and 5 & Local around nearly uniform logits; may yield negative approximate weights. \\
\badge{Purple}{ANSATZ} & Equation 7 & A finite affine state cannot represent generic softmax dependence on every query exactly. \\
\badge{Orange}{HEURISTIC} & Equation 17 & Direction follows under a norm constraint; magnitude assumes comparable query/key scales. \\
\badge{WechatGreen}{STRUCTURE} & Equations 18--19 & Retain--erase--write forms match after centering/rescaling; learned models are not thereby equal. \\
\badge{Red}{NOT CLAIMED} & Performance & The article contains no experiment, benchmark, accuracy guarantee, or universal replacement claim. \\
\bottomrule
\end{longtable}

Additional boundaries matter in practice:

\begin{itemize}
  \setlength{\itemsep}{0.35em}
  \item Repeated local approximation errors may accumulate through the scan.
  \item The minimax-inspired anchor does not control the replacement error for
        every possible future query.
  \item The derivation assumes causal dot-product attention. Extra biases and
        masks need separate analysis, as do positional terms and nonlinear
        feature maps.
  \item The deterministic $1/t$ write rate becomes small on long prefixes;
        learned GDN gates need not follow this schedule.
  \item Centered keys can have varying norms, so practical implementations may
        need normalization or learned scaling not justified by this derivation.
\end{itemize}

\section{Research ideas suggested by the mechanism}

The following are testable follow-ups, not claims made by the source article.

\subsection{Use logit dispersion as an approximation-confidence signal}

The Taylor error grows with centered-logit spread. A model could estimate
$\operatorname{Var}_i(\vect q_t\trans\vect k_i)$ or a recurrent proxy and use
it to interpolate between the cheap Delta state and a local attention fallback.
This directly tests whether the mathematical validity region predicts model
error.

\subsection{Initialize learned gates from the harmonic schedule}

The derivation suggests $g_t=1/t$ and $\lambda_t=1-1/t$ as an analytic starting
point. A learned GDN could be initialized near this schedule, then trained to
depart from it when the data calls for faster writing or forgetting. The useful
comparison is initialization and stability, not an assumption that the fixed
schedule is optimal.

\subsection{Ablate centering separately from the Delta Rule}

The derived key is $\vect k_t-\bar{\vect k}_t$ and the value target is
$\vect v_t-\bar{\vect v}_{t-1}$. A controlled experiment should compare
uncentered, key-centered, value-centered, and doubly centered updates while
holding parameter count and gate parameterization fixed.

\subsection{Measure diagonal approximation error over time}

Equation~\eqref{eq:factor} provides an interpretable diagnostic: compare exact
diagonal weights, their Taylor approximations, and the product-induced old
weights. This can reveal whether the local-factor strategy really outperforms a
single global expansion in trained representations, and where errors compound.

\subsection{Benchmark quality and systems behavior separately}

Any evaluation should report approximation error, language-model quality,
long-context retrieval, recurrent-state memory, training throughput, and
single-token decoding latency. Linear asymptotics and hardware speed are
different hypotheses.

\begin{greenbox}{FINAL TAKEAWAY}
  Softmax contains an exact Delta-shaped update when the query is fixed. The
  article linearizes only the newest attention gate, exposes the query term that
  prevents a shared state, and anchors that term to the centered newest key.
  The resulting matrix recurrence retains old memory, erases its prediction in
  one key direction, and writes the residual there. That is why Gated
  DeltaNet's mechanism appears---approximately and structurally, not by exact
  model equivalence.
\end{greenbox}

\appendix
\section{Dimension check for the final update}

For $\vect A_{t-1}\in\R^{d_v\times d_k}$,
$\vect x_t\in\R^{d_k}$, and $\vect y_t\in\R^{d_v}$:

\[
  \underbrace{\vect A_{t-1}\vect x_t}_{d_v\text{-vector}},
  \qquad
  \underbrace{\vect y_t-\vect A_{t-1}\vect x_t}_{d_v\text{-vector}},
  \qquad
  \underbrace{(\vect y_t-\vect A_{t-1}\vect x_t)\vect x_t\trans}_{d_v\times d_k\text{ matrix}}.
\]

Every term in Equation~\eqref{eq:deriveddelta} therefore has the same
$d_v\times d_k$ shape. Reading with $\vect q_t\in\R^{d_k}$ produces a
$d_v$-vector, which can be added to $\bar{\vect v}_t$.

\section{Source-equation map}

\begin{tabularx}{\textwidth}{@{}P{18mm}P{43mm}Y@{}}
\toprule
\textbf{Source} & \textbf{Role here} & \textbf{Status} \\
\midrule
1--2 & Attention definition and global first-order VaLA baseline & Definition, then local approximation \\
3--4 & Prefix split and prediction-error recurrence & Exact for a fixed query \\
5--6 & Diagonal Taylor step and survival factorization & Approximation, then exact identity \\
7--10 & Affine state and quadratic-query obstruction & Modeling ansatz and inherited approximation \\
11--14 & Dropped-query shortcut and return to VaLA & Heuristic followed by exact telescoping algebra \\
15--17 & Replacement error and representative direction & Exact error expression followed by heuristic choice \\
18--20 & GDN-shaped update, canonical comparison, and readout & Derived approximation and structural comparison \\
\bottomrule
\end{tabularx}

\bibliographystyle{plain}
\bibliography{refs}

\end{document}
